\chapter{Principio dei lavori virtuali}
\begin{definition}{}[Lavoro di una forza]
Sia \(\underline{f}(\underline{x},\underline{v},t)\) una forza che agisce nel punto materiale di posizione \(\underline{x}(t)\), e velocità \(\underline{v}(t)\). Definiamo il lavoro della forza lungo la traiettoria \(\underline{x}(t)\) come
\[
L=\int_{\underline{x}} \underline{f}\cdot d\underline{x}
\]

\begin{center}
\begin{tikzpicture}[scale=1.1, >=Latex]

% assi
\draw[->, line width=0.8pt] (0,0) -- (-2.5,-2.7);
\draw[->, line width=0.8pt] (0,0) -- (3.2,0);
\draw[->, line width=0.8pt] (0,0) -- (0,2.3);

\node at (-2.65,-2.85) {\(x\)};
\node at (3.25,-0.45) {\(y\)};
\node at (0.35,2.35) {\(z\)};

% traiettoria
\draw[line width=0.8pt]
  (0.25,-1.2)
  .. controls (0.45,-0.2) and (0.8,0.4) ..
  (1.6,0.7)
  .. controls (2.4,1.0) and (2.9,1.7) ..
  (3.35,2.1);

% estremi e punto materiale
\fill (0.25,-1.2) circle (1.7pt);
\fill (3.35,2.1) circle (1.7pt);
\fill (1.7,0.62) circle (1.7pt);

\node at (1.05,0.95) {\(\underline{x}(t)\)};

% forza
\draw[->, line width=0.8pt] (1.7,0.62) -- (2.85,0.15);
\node at (2.75,0.7) {\(\underline{f}(\underline{x},\underline{v},t)\)};

\end{tikzpicture}
\end{center}
\end{definition}
Ricordiamo che ad ogni vincolo è associata una reazione vincolare che assicura il soddisfacimento del vincolo stesso.

\begin{definition}{}[Vincoli perfetti]
Un vincolo perfetto è un vincolo per il quale la reazione vincolare non compie lavoro, perché è \(\perp\) (localmente) alla traiettoria. Matematicamente un vincolo perfetto è tale che
\[
\underline{\Psi}\cdot \delta \underline{x}=0
\qquad
\forall\, \delta \underline{x}\ \text{compatibile con il vincolo}
\]
\end{definition}


Risulta necessario definire con chiarezza cosa siano i \(\delta \underline{x}\) e cosa voglia dire che questi siano ``compatibili con il vincolo''.
Lo facciamo nel caso in cui i vincoli siano regolari ma è possibile farlo anche in caso contrario.
Per fare ciò in modo matematicamente rigoroso abbiamo bisogno di strumenti di geometria differenziale avanzata e calcolo tensoriale che vanno al di là degli scopi del corso. Presenteremo quindi una trattazione meno rigorosa e più intuitiva basandoci su concetti di geometria differenziale di base rifacendoci soprattutto al caso noto delle superfici regolari.
\inizioapprof
I \(\delta \underline{x}\) vengono detti spostamenti virtuali. Supponiamo che un punto

materiale \((\underline{x}(t),m)\) sia soggetto al vincolo olonomo \(\overline{\Phi}(\underline{x},t)=0\).
Per ogni istante \(t\) definiamo la sottovarietà
\[
M_t=\{P\in\mathbb{R}^3:\overline{\Phi}(P,t)=0\}
\]

ovvero \(M_t=\overline{\Phi}^{-1}(0,t)\). Nel caso semplice in cui \(\underline{x}(t)\in\mathbb{R}^3\) questa può essere una superficie o una curva. In un caso più generale questa è una sotto-varietà dello spazio delle configurazioni \(Q\).
Se il vincolo è regolare, ovvero se \(d\overline{\Phi}(\underline{x},t)\) ha rango massimo \(\forall x\in\mathbb{R}^3\), allora \(M_t\) è una sottovarietà regolare. Risulta quindi definito in ogni punto \(\underline{x}\in M_t\) lo spazio tangente \(T_xM_t=\Ker(d\overline{\Phi}(\underline{x},t))\) a \(M_t\).
Per capire più facilmente questa condizione possiamo pensare al caso di superficie regolare ottenuta tramite controimmagine di un valore regolare: sia \(F:\mathbb{R}^3\to\mathbb{R}\) differenziabile e sia \(S=F^{-1}(0)\) con \(dF_p\) di rango massimo \(\forall p\in S\). In questo caso possiamo usare l'identificazione del funzionale \(dF_p:\mathbb{R}^3\to\mathbb{R}\) con il gradiente \(\nabla F_p\in\mathbb{R}^3\) data da \(dF_p(v)=\nabla F_p^Tv\ \ \forall v\in\mathbb{R}^3\).
Dire che \(dF_p\) ha rango massimo equivale a dire \(\nabla F_p\neq 0\) infatti \(\mathrm{rg}(dF_p)=\mathrm{rg}(L_{\nabla F_p})=0\) sse \(\nabla F_p=0\).
Allora lo spazio tangente non è altro che \(T_pS=\Ker(dF_p)=\{v\in\mathbb{R}: \nabla F_p^Tv=0\}\).\\
Gli spostamenti virtuali \(\delta \underline{x}\) sono in ogni istante fissato \(t\) gli elementi di \(T_{\underline{x}(t)}M_t\), ovvero sono le direzioni percorribili compatibili con la varietà \(M_t\) a \(t\) fissato.
È importante non confondere gli spostamenti virtuali \(\delta \underline{x}\) con gli spostamenti infinitesimi reali \(d\underline{x}=\underline{\dot{x}}(t)\,dt\).

I \(\delta \underline{x}\) sono tutti i possibili spostamenti percorribili in un tempo infinitesimo, sono una linearizzazione istantanea dei movimenti ammessi ma non è detto che questi vengano effettuati dalla traiettoria \(\underline{x}(t)\) (anzi al massimo uno di loro potrebbe essere effettuato ma non è detto che lo sia).

Dall'altra parte \(d\underline{x}=\underline{\dot{x}}(t)\,dt\) è effettivamente lo spostamento effettuato in un tempo infinitesimo \(dt\) dopo l'istante \(t\).\\
La differenza è netta tra i due nel caso di un vincolo dipendente dal tempo \(\overline{\Phi}(\underline{x},t)=0\). Immaginiamo due istanti di tempo \(t_1,t_2\) infinitesimamente vicini ai quali corrispondono le varietà \(M_{t_1},M_{t_2}\) con \(\underline{x}(t_1)\in M_{t_1}\), \(\underline{x}(t_2)\in M_{t_2}\).
Gli spostamenti virtuali stanno in \(T_{\underline{x}(t_1)}M_{t_1}\) e in \(T_{\underline{x}(t_2)}M_{t_2}\) che sono due spazi ben distinti. Lo spostamento reale infinitesimo invece è
\[
d\underline{x}=\underline{x}(t_2)-\underline{x}(t_1)\qquad (\text{abbiamo supposto } t_1-t_2\sim 0)
\]

che non sta in nessuno dei 2 tangenti.
Per capire meglio questo concetto pensiamo al caso del vincolo
\[
\overline{\Phi}(x,y,t)=(x-x_c(t))^2+(y-y_c(t))^2-R(t)^2=0
\]
Ad ogni istante \(t\) il punto è vincolato a muoversi su una circonferenza di raggio e centro variabile. Disegniamo spostamenti virtuali e spostamento reale per 2 istanti \(t_1,t_2\) molto vicini. \(C_1=(x_c(t_1),y_c(t_1))\), \(C_2=(x_c(t_2),y_c(t_2))\)
\begin{center}
\begin{tikzpicture}[scale=1,>=Latex]

% assi cartesiani
\draw[->, line width=1pt] (0.8,0.8) -- (0.8,6.2);
\draw[->, line width=1pt] (0.8,0.8) -- (9.4,0.8);
\node at (0.45,5.95) {\(y\)};
\node at (9.15,0.45) {\(x\)};

% circonferenze
\draw[line width=0.9pt] (3.0,2.2) circle (1.2);
\draw[line width=0.9pt] (7.4,4.2) circle (1.7);

% centri
\fill (3.0,2.2) circle (2pt);
\fill (7.4,4.2) circle (2pt);
\node at (2.75,1.75) {\(C_1\)};
\node at (7.15,3.75) {\(C_2\)};

% punti x(t1), x(t2)
\fill[purple] (3.8485,3.0485) circle (2.2pt);
\fill[purple] (6.3070,2.8978) circle (2.2pt);

\node[purple] at (4.55,3.35) {\(\underline{x}(t_1)\)};
\node[purple] at (6.95,2.45) {\(\underline{x}(t_2)\)};

% raggi
\draw[line width=0.9pt] (3.0,2.2) -- (3.8485,3.0485);
\draw[line width=0.9pt] (7.4,4.2) -- (6.3070,2.8978);

\node at (3.15,3.00) {\(R(t_1)\)};
\node at (6.45,3.95) {\(R(t_2)\)};

% tangente in x(t1)
\draw[line width=1pt, color=green!60!black] (3.32,3.58) -- (4.38,2.52);
\node[color=green!60!black] at (2.55,3.95) {\(T_{\underline{x}(t_1)}M\)};

% tangente in x(t2)
\draw[line width=1pt, color=red!80!black] (5.62,3.47) -- (6.99,2.33);
\node[color=red!80!black] at (8.25,2.55) {\(T_{\underline{x}(t_2)}M\)};

% vettore dx
\draw[->, line width=1pt, color=orange!90!black] (3.8485,3.0485) -- (6.3070,2.8978);
\node[color=orange!90!black] at (5.05,3.22) {\(\underline{dx}\)};

\end{tikzpicture}
\end{center}

Nel caso di vincolo costante si ha che \(M_t=M\ \forall t\) e che \(\underline{x}\in M\), \(d\underline{x}\in T_{\underline{x}}M\). Anche in questo caso però lo spostamento reale è un'entità distinta dagli spostamenti virtuali perché è l'unico che viene effettivamente effettuato.

\begin{oss}{}
Il concetto di spostamento virtuale si può generalizzare a un punto materiale soggetto a \(m\) vincoli olonomi differenziabili
\[
\overline{\Phi}_i(\underline{x},t)=0
\qquad
\forall i=1,\ldots,m
\]
in tal caso lo spazio dei \(\delta \underline{x}\) è \(T_{\underline{x}}M\) con
\[
M=\bigcap_{i=1}^{m}\Ker(d\overline{\Phi}_i)\Rightarrow \forall\,\delta \underline{x}\in T_{\underline{x}}M \quad d\overline{\Phi}_i\,\delta \underline{x}=0
\]

Allo stesso modo possiamo definire il concetto di spostamento virtuale per un sistema di \(n\) punti \(\underline{x}_1,\ldots,\underline{x}_n\) soggetto a \(m\) vincoli olonomi perfetti \(\overline{\Phi}_1(\underline{x}_1,\ldots,\underline{x}_n,t)\ldots \overline{\Phi}_m(\underline{x}_1,\ldots,\underline{x}_n,t)\).

Indichiamo una configurazione del sistema come
\[
\underline{x}=(\underline{x}_1,\ldots,\underline{x}_n)\in\mathbb{R}^{3n}.
\]
Imponendo i vincoli otteniamo
\[
M_t=\left\{\underline{x}\in\mathbb{R}^{3n}:\overline{\Phi}_1(\underline{x},t)=\ldots=\overline{\Phi}_m(\underline{x},t)=0\right\}
\]

Se i vincoli sono indipendenti e regolari su \(M_t\), ovvero il loro differenziale ha rango massimo su \(M_t\), quest'ultima è una sotto varietà regolare di \(\mathbb{R}^{3n}\) di dimensione \(r=3n-m\).

Uno spostamento virtuale \(\delta \underline{x}\) del sistema nella configurazione \(\underline{x}\in M_t\) è un vettore del tangente in \(\underline{x}\) a \(M_t\)
\[
\delta \underline{x}=(\delta \underline{x}_1,\ldots,\delta \underline{x}_n)\in T_{\underline{x}}M_t
\]

Importante notare che le \(\delta \underline{x}_i\) non sono arbitrarie proprio perché devono garantire \(\delta \underline{x}\in T_{\underline{x}}M_t\).

Si ha che
\[
T_{\underline{x}}M_t=\bigcap_{\nu=1}^{m}\Ker(d\overline{\Phi}_{\nu\,\underline{x}})=\left\{\delta \underline{x}:d\overline{\Phi}_{\nu\,\underline{x}}(\delta \underline{x})=0 \quad \forall \nu=1,\ldots,m\right\}
\]

Scrivendo in componenti
\[
T_{\underline{x}}M_t=\left\{\delta \underline{x}:\nabla \overline{\Phi}_\nu(\underline{x},t)^T\delta \underline{x}=0 \quad \forall \nu=1,\ldots,m\right\}=\mathrm{span}^{\perp}\left\{\nabla \overline{\Phi}_1(\underline{x},t),\ldots,\nabla \overline{\Phi}_m(\underline{x},t)\right\}
\]

ovvero è lo spazio ortogonale ai gradienti di tutti i vincoli in \(\underline{x}\).
\end{oss}

\begin{theorem}{}[Principio dei lavori virtuali, ipotesi semplificate, vincoli dipendenti solo da un punto]
Consideriamo un sistema meccanico formato dai punti \(\underline{x}_1,\ldots,\underline{x}_n\). Il punto \(\underline{x}_i\) è soggetto ai vincoli perfetti regolari
\[
\overline{\Phi}_{i,j}(\underline{x}_i)=0
\qquad
\forall j=1,\ldots,m_i.
\]
In ogni punto \(\underline{x}_i\) è esercitata una forza risultante
\[
\underline{f}_i=\underline{f}_i^{\mathrm{att.}}+\underline{f}_i^{\mathrm{reat.}}
\]
All'equilibrio vale la relazione
\[
\sum_{i=1}^{n}\underline{f}_i\cdot \delta \underline{x}_i=0
\qquad
\forall (\delta \underline{x}_1,\ldots,\delta \underline{x}_n)\in T_{\underline{x}_1}M_1\times\cdots\times T_{\underline{x}_n}M_n
\]
dove \(M_i\) è la varietà generata dai vincoli \(\{\overline{\Phi}_{i,j}(\underline{x}_i)\}_{j=1}^{m_i}\) e gli \(\delta \underline{x}_i\) sono gli spostamenti virtuali compatibili con i vincoli.
\end{theorem}

\begin{oss}{}
Similmente al concetto dei moltiplicatori di Lagrange stiamo dicendo che \(\forall i\) se \(\overline{\Phi}_{i,j}(\underline{x}_i)=0\) \(j=1,\ldots,m_i\): sono i vincoli perfetti imposti sul punto \(\underline{x}_i\) allora
\[
\underline{f}_i\in \mathrm{Span}\left\{\nabla \overline{\Phi}_{i,1}(\underline{x}_i),\ldots,\nabla \overline{\Phi}_{i,m_i}\right\}
\Longleftrightarrow
\underline{f}_i=\lambda_1\nabla \overline{\Phi}_{i,1}(\underline{x}_i)+\cdots+\lambda_{m_i}\nabla \overline{\Phi}_{i,m_i}
\]
\end{oss}
\begin{proof}
Sia \(M_i=\{\underline{x}_i\in\mathbb{R}^3:\overline{\Phi}_{i,1}(\underline{x}_i)=\ldots=\overline{\Phi}_{i,m_i}(\underline{x}_i)=0\}\) e sia \((\underline{x}_1,\ldots,\underline{x}_n)\in M_1\times\ldots\times M_n\).

Per ogni punto \(\underline{x}_i\) vale \(\underline{f}_i=\underline{f}_i^{\mathrm{att}}+\underline{f}_i^{\mathrm{reat}}\). Poiché i vincoli sono perfetti per ogni scelta di spostamenti virtuali ammessi vale
\[
\sum_{i=1}^{n}\underline{f}_i^{\mathrm{reat}}\cdot \delta \underline{x}_i=0
\qquad
(\delta \underline{x}_1,\ldots,\delta \underline{x}_n)\in T_{\underline{x}_1}M_1\times\ldots\times T_{\underline{x}_n}M_n
\]

Quindi la condizione sui lavori virtuali equivale a una sulle sole forze attive
\[
\sum_{i=1}^{n}\underline{f}_i\cdot \delta \underline{x}_i=0
\Longleftrightarrow
\sum_{i=1}^{n}\underline{f}_i^{\mathrm{att}}\cdot \delta \underline{x}_i=0
\]

Basta quindi dimostrare il teorema per la seconda condizione

\((\Rightarrow)\) Se siamo in una configurazione di equilibrio allora \(\forall i=1,\ldots,n\)
\[
\underline{f}_i^{\mathrm{att}}+\underline{f}_i^{\mathrm{reat}}=0
\]

Moltiplicando scalarmente per un qualsiasi \((\delta \underline{x}_1,\ldots,\delta \underline{x}_n)\) si ottiene
\[
\sum_{i=1}^{n}\underline{f}_i^{\mathrm{att}}\cdot \delta \underline{x}_i=
\sum_{i=1}^{n}\left(\underline{f}_i^{\mathrm{att}}+\underline{f}_i^{\mathrm{reat}}\right)\cdot \delta \underline{x}_i=0
\]

\((\Leftarrow)\) Supponiamo ora che
\[
\sum_{i=1}^{n}\underline{f}_i^{\mathrm{att}}\cdot \delta \underline{x}_i=0
\qquad
(\delta \underline{x}_1,\ldots,\delta \underline{x}_n)\in T_{\underline{x}_1}M_1\times\ldots\times T_{\underline{x}_n}M_n
\]

In particolare se scelgo un \(i\in\{1,\ldots,n\}\) deve essere vero per una configurazione del tipo \((0,\ldots,0,\delta \underline{x}_i,0,\ldots,0)\) con \(\delta \underline{x}_i\in T_{\underline{x}_i}M_i\) arbitrario. La condizione diventa
\[
\underline{f}_i^{\mathrm{att}}\cdot \delta \underline{x}_i=0
\qquad
\forall \delta \underline{x}_i\in T_{\underline{x}_i}M_i
\]

Quindi \(\underline{f}_i^{\mathrm{att}}\in (T_{\underline{x}_i}M_i)^{\perp}=\mathrm{span}\{\nabla \overline{\Phi}_{i,1}(\underline{x}_i),\ldots,\nabla \overline{\Phi}_{i,m_i}(\underline{x}_i)\}\Longleftrightarrow \exists \lambda_{i,1},\ldots,\lambda_{i,m_i} \in \mathbb{R}\) tali che
\[
\underline{f}_i^{\mathrm{att}}=\sum_{j=1}^{m_i}\lambda_{i,j}\nabla \overline{\Phi}_{i,j}(\underline{x}_i)
\]

Definiamo allora \(\underline{f}_i^{\mathrm{reat}}=-\sum_{j=1}^{m_i}\lambda_{i,j}\nabla \overline{\Phi}_{i,j}(\underline{x}_i)\). Ovviamente si ha \(\underline{f}_i^{\mathrm{att}}+\underline{f}_i^{\mathrm{reat}}=0\).

Abbiamo trovato una reazione vincolare compatibile con i vincoli che determina una configurazione di equilibrio valida per \(\underline{x}_i\). Dal momento che la \(i\) era arbitraria ho trovato una configurazione statica ammissibile per il sistema.
\end{proof}
\begin{theorem}{}[Principio dei lavori virtuali, vincoli regolari]
Consideriamo un sistema meccanico formato dai punti \(\underline{x}_1,\ldots,\underline{x}_n\) soggetto a \(m\) vincoli olonomi perfetti regolari
\[
\overline{\Phi}_1(\underline{x}_1,\ldots,\underline{x}_n,t)=\ldots=\overline{\Phi}_m(\underline{x}_1,\ldots,\underline{x}_n,t)=0.
\]
Supponiamo inoltre che in ogni punto agisca una risultante delle forze attive puramente posizionali \(\underline{f}_i^{\mathrm{att}}\). Allora siamo in una configurazione di equilibrio sse
\[
\sum_{i=1}^{n}\underline{f}_i^{\mathrm{att}}\cdot \delta \underline{x}_i=0
\qquad
\forall\, \delta \underline{x}=(\delta \underline{x}_1,\ldots,\delta \underline{x}_n)\in T_{\underline{x}}M_t
\]
dove \(M_t\) è la sottovarietà generata dai vincoli.
\end{theorem}
\begin{oss}{}
    Anche il principio dei lavori virtuali è valido nel caso in cui i vincoli siano solo olonomi perfetti ma non regolari. In questo caso però si perde la trattazione data dalla geometria differenziale e si deve procedere in altro modo.
\end{oss}
\begin{proof}
Indichiamo con \(\underline{R}=(\underline{R}_1,\ldots,\underline{R}_n)\) il sistema delle reazioni vincolari, allora dal momento che i vincoli sono perfetti vale
\[
\sum_{i=1}^{n}\underline{R}_i\cdot \delta \underline{x}_i=0
\qquad
\forall\, \delta \underline{x}=(\delta \underline{x}_1,\ldots,\delta \underline{x}_n)\in T_{\underline{x}}M_t
\]
ovvero \(\underline{R}\in T_{\underline{x}}M_t^\perp\).

\((\Rightarrow)\) Supponiamo \(\underline{x}\) di equilibrio, per definizione esistono le reazioni vincolari \(\underline{R}_1,\ldots,\underline{R}_n\) tali che
\[
\underline{f}_i^{\mathrm{att}}+\underline{R}_i=0
\qquad
\forall i=1,\ldots,n
\]
Sia \(\delta \underline{x}\) uno spostamento virtuale arbitrario. Allora
\[
0=\sum_{i=1}^{n}\left(\underline{f}_i^{\mathrm{att}}+\underline{R}_i\right)\cdot \delta \underline{x}_i
=\sum_{i=1}^{n}\underline{f}_i^{\mathrm{att}}\cdot \delta \underline{x}_i
\]

\((\Leftarrow)\) Supponiamo ora che
\[
\sum_{i=1}^{n}\underline{f}_i^{\mathrm{att}}\cdot \delta \underline{x}_i=0
\qquad
\forall\, \delta \underline{x}=(\delta \underline{x}_1,\ldots,\delta \underline{x}_n)\in T_{\underline{x}}M_t
\]
Introduciamo il vettore globale delle forze attive
\[
\underline{A}=(\underline{f}_1^{\mathrm{att}},\ldots,\underline{f}_n^{\mathrm{att}})\in\mathbb{R}^{3n},
\]
la condizione si riscrive come
\[
\underline{A}\cdot \delta \underline{x}=0
\qquad
\forall\, \delta \underline{x}\in T_{\underline{x}}M_t
\Longleftrightarrow
\underline{A}\in T_{\underline{x}}M_t^\perp
\]
Siccome i vincoli sono regolari
\[
T_{\underline{x}}M_t^\perp=\mathrm{Span}\left\{\nabla \overline{\Phi}_{1\underline{x}},\ldots,\nabla \overline{\Phi}_{m\underline{x}}\right\},
\]
perciò esistono $\lambda_1, \dots, \lambda_m \in \mathbb{R}$ tali che
\[
\underline{A}=\sum_{j=1}^{m}\lambda_j\,\nabla \overline{\Phi}_{j\underline{x}}
\]
Definiamo ora \(\underline{R}=-\underline{A}\). Per definizione \(\underline{R}\in T_{\underline{x}}M_t^\perp\) ed è quindi una reazione vincolare ammissibile. Se riscriviamo la condizione \(\underline{A}+\underline{R}=0\) per blocchi otteniamo
\[
(\underline{f}_1^{\mathrm{att}},\ldots,\underline{f}_n^{\mathrm{att}})+(\underline{R}_1,\ldots,\underline{R}_n)=0
\]
che scritta componente per componente è
\[
\underline{f}_i^{\mathrm{att}}+\underline{R}_i=0
\qquad
\forall i=1,\ldots,n
\]
ovvero la condizione di equilibrio.

Abbiamo quindi dimostrato che il principio dei lavori virtuali è condizione necessaria e sufficiente per l'equilibrio.
\end{proof}
\fineapprof
Consideriamo un sistema meccanico con \(n\) punti, \(3n\) gradi di libertà, \(m\) vincoli olonomi perfetti regolari.

Poiché i vincoli sono anche olonomi possiamo determinare \(q_1,\ldots,q_{3n-m}\) coordinate libere.
Le posizioni \(\underline{x}_i(t)\) sono univocamente determinate dalle coordinate libere
\[
\underline{x}_i(q_1(t),\ldots,q_{3n-m}(t))
\]

Possiamo considerare gli spostamenti virtuali nello spazio delle variabili libere.
Localmente il differenziale di \(\underline{x}_i(q_1,\ldots,q_{3n-m})\) manda un \(\delta \underline{q}=(\delta q_1,\ldots,\delta q_{3n-m})\) in un \(\delta \underline{x}_i\), ovvero
\[
\delta \underline{x}_i=d\underline{x}_{i,\underline{q}}(\delta \underline{q})
\]

In coordinate
\[
\delta \underline{x}_i=\sum_{j=1}^{3n-m}\frac{\partial \underline{x}_i}{\partial q_j}(\underline{q})\,\delta q_j
\]

\begin{example}{}
Consideriamo un punto che si muove su una circonferenza. \(i=1\) \(q=\theta\).
\[
\underline{x}_1=(R\cos\theta,R\sin\theta)
\]

\[
\delta \underline{x}=\frac{\partial \underline{x}}{\partial q}\cdot \delta q
=\frac{\partial \underline{x}}{\partial \theta}\cdot \delta \theta
=\frac{\partial}{\partial \theta}
\begin{pmatrix}
R\cos\theta\\
R\sin\theta
\end{pmatrix}
\delta \theta
=
\begin{pmatrix}
-R\sin\theta\,\delta \theta\\
R\cos\theta\,\delta \theta
\end{pmatrix}
\]
\end{example}

Se metto questa rappresentazione dello spostamento virtuale in funzione delle coordinate libere nel principio dei lavori virtuali ottengo
\[
\sum_{i=1}^{n}\underline{f}_i\cdot \delta \underline{x}_i
=
\sum_{i=1}^{n}\underline{f}_i^{\mathrm{att}}\cdot \delta \underline{x}_i
=
\sum_{i=1}^{n}\underline{f}_i^{\mathrm{att}}\cdot \sum_{j=1}^{3n-m}\frac{\partial \underline{x}_i}{\partial q_j}\,\delta q_j
=
\sum_{j=1}^{3n-m}\left(\sum_{i=1}^{n}\underline{f}_i\cdot \frac{\partial \underline{x}_i}{\partial q_j}\right)\delta q_j
\]

Chiamiamo
\[
\sum_{i=1}^{n}\underline{f}_i\cdot \frac{\partial \underline{x}_i}{\partial q_j}=Q_j
\qquad \Rightarrow \qquad
\sum_{i=1}^{n}\underline{f}_i\cdot \delta \underline{x}_i
=
\sum_{j=1}^{3n-m}Q_j\,\delta q_j
\]

Possiamo dare un'altra dimostrazione del principio dei lavori virtuali basato sulla parametrizzazione locale della varietà attraverso le coordinate libere.
\begin{proof}
Alternativamente si può procedere in coordinate libere. Si ha
\[
\sum_{i=1}^{n} \underline{f}_i^{\mathrm{att}}\cdot \delta \underline{x}_i
=
\sum_{i=1}^{n} \underline{f}_i^{\mathrm{att}}\cdot
\sum_{j=1}^{3n-m}
\frac{\partial \underline{x}_i}{\partial q_j}\,\delta q_j
=
\sum_{j=1}^{3n-m}
\left(
\underbrace{
\sum_{i=1}^{n}
\underline{f}_i^{\mathrm{att}}\cdot
\frac{\partial \underline{x}_i}{\partial q_j}
}_{Q_j}
\right)
\delta q_j.
\]

Dal momento che
\[
\underline{Q}\cdot
\begin{pmatrix}
\delta q_1\\
\vdots\\
\delta q_{3n-m}
\end{pmatrix}
=0
\qquad
\forall\,(\delta q_1,\dots,\delta q_{3n-m}),
\]
segue che
\[
\underline{Q}=\underline{0}
\qquad \Longrightarrow \qquad
Q_j=0
\quad
\forall j=1,\dots,3n-m.
\]

Allora la risultante delle forze attive è perpendicolare allo spazio delle direzioni ammesse dai vincoli e arriviamo alla stessa conclusione di prima.
\end{proof}

\begin{example}{}
  Studiamo la statica del seguente sistema.
  \begin{center}
    \begin{tikzpicture}[
    x=1cm,
    y=1cm,
    >=Latex,
    line cap=round,
    line join=round,
    every node/.style={font=\small}
]

% Parametri
\def\R{1.55}

% Coordinate
\coordinate (O) at (0,0);
\coordinate (A) at (0,3.45);
\coordinate (G) at (7.20,\R);
\coordinate (C) at (7.20,0);
\coordinate (Top) at ($(G)+(0,\R)$);
\coordinate (B) at ($(A)!0.42!(G)$);

% Assi
\draw[black, line width=1.4pt, ->] (0,0) -- (10.2,0) node[below right] {$x$};
\draw[black, line width=1.4pt, ->] (0,0) -- (0,4.2) node[above left] {$y$};

% Parete verticale
\draw[black, line width=2pt] (0,0) -- (0,4.05);

% Origine
\node[below left] at (O) {$O$};

% Riferimento centrato in O
\draw[orange!85!black, line width=1.4pt, ->] (O) -- ++(0.95,0);
\draw[orange!85!black, line width=1.4pt, ->] (O) -- ++(0,0.95);
\node[orange!85!black, below] at ($(O)+(0.48,0)$) {$\underline{\imath}$};
\node[orange!85!black, left] at ($(O)+(0,0.48)$) {$\underline{\jmath}$};

\draw[orange!85!black, line width=1.2pt] (-0.45,0.18) circle (0.18);
\fill[orange!85!black] (-0.45,0.18) circle (0.03);
\node[orange!85!black, left] at (-0.62,0.62) {$\underline{k}$};

% Sbarrette verticali attaccate all'asse y
\draw[black, line width=1.6pt] (0,3.14) -- (0,3.56);
\draw[black, line width=1.6pt] (0.12,3.14) -- (0.12,3.56);

% Punto A
\node[left] at (A) {$A$};

% Asta
\draw[black, line width=1.8pt] (A) -- (G);

% Punto B
\fill[red] (B) circle (2.8pt);
\node[above] at (B) {$B$};

% Etichetta asta
\node[above right] at ($(B)!0.52!(G)$) {$l,m$};

% Angolo theta
\draw[black, line width=1.1pt]
($(A)+(0,-0.78)$) arc[start angle=-90,end angle=-14.5,radius=0.78];
\node at ($(A)+(0.72,-0.87)$) {$\theta$};

% Freccia verde in A
\draw[green!55!black, line width=1.4pt, ->] ($(A)$) -- ++(1.80,0);
\node[green!55!black, above] at ($(A)+(1.25,0.20)$) {$\underline{\psi}$};

% Molla attaccata all'asse y
\draw[black, line width=1.4pt, smooth, samples=240, domain=0:7.02]
plot (\x,{1.55 + 0.10*sin(900*\x)});

% Ruota
\draw[black, line width=1.8pt] (G) circle (\R);

% Punto C
\node[below] at (C) {$C$};

% Centro G
\filldraw[white, draw=black, line width=1.2pt] (G) circle (0.09);
\node[above left] at ($(G)+(-0.10,0.08)$) {$G$};

% Etichetta M,R
\node[above right] at ($(Top)+(0.95,-0.05)$) {$M,R$};

% Raggi tratteggiati
\draw[black, dashed, line width=1.1pt] (Top) -- (G);
\draw[black, dashed, line width=1.1pt] (G) -- ++(0.95,1.00);

% Angolo phi
\draw[black, line width=1.0pt, ->]
($(G)+(0,0.58)$) arc[start angle=90,end angle=46,radius=0.58];
\node at ($(G)+(0.47,0.92)$) {$\varphi$};

% Forza mg in B
\draw[red!85!black, line width=1.5pt, ->] (B) -- ++(0,-1.15);
\node[red!85!black, left] at ($(B)+(0.10,-0.70)$) {$m\underline{g}$};

% Forza Mg in G
\draw[red!85!black, line width=1.5pt, ->] (G) -- ++(0,-1.05);
\node[red!85!black, right] at ($(G)+(0.52,-0.62)$) {$M\underline{g}$};

% Forza elastica
\draw[red!85!black, line width=1.5pt, ->] (G) -- ++(-1.25,0);
\node[red!85!black, below] at ($(G)+(-0.82,-0.12)$) {$-k\,\underline{x}_G$};

% Assi verdi vicino a C
\draw[green!55!black, line width=1.4pt, ->] ($(C)+(0.18,0)$) -- ++(1.90,0);
\draw[green!55!black, line width=1.4pt, ->] ($(C)+(0.18,0)$) -- ++(0,0.95);
\node[green!55!black, right] at ($(C)+(2.10,0)$) {$\underline{\xi}_x$};
\node[green!55!black, right] at ($(C)+(0.22,0.58)$) {$\underline{\xi}_y$};

\end{tikzpicture}
  \end{center}
Imposto un riferimento fisso
\[
\{O,\hat{\imath},\hat{\jmath},\hat{k}\}.
\]
Il sistema è soggetto ai vincoli
\[
\left\{
\begin{array}{l}
x_G^{\mathrm{disco}}=x_G^{\mathrm{asta}},\\[2mm]
y_G^{\mathrm{disco}}=y_G^{\mathrm{asta}}.
\end{array}
\right.
\]

Come coordinata libera posso scegliere $\theta$, $\psi$ o $x_G$.
Supponendo che per $\psi=0$ il corpo sia in $x_G=0$, si ha
\[
l\sin\theta = R\psi = x_G.
\]

Scrivo le equazioni della statica per i due corpi:
\[
\left\{
\begin{array}{l}
\underline{\psi}
+
m\underline{g}
+
M\underline{g}
-
k\underline{x}_G
+
\underline{\xi}_x
+
\underline{\xi}_y
=
\underline{0},
\\[2mm]
\left(\underline{x}_C-\underline{x}_G\right)\wedge \underline{\xi}_x
=
\underline{0},
\\[2mm]
\left(\underline{x}_B-\underline{x}_G\right)\wedge m\underline{g}
+
\left(\underline{x}_A-\underline{x}_G\right)\wedge \underline{\psi}
=
\underline{0}.
\end{array}
\right.
\]

Dove nella prima ho accorpato la prima legge cardinale per i due corpi
attraverso il principio di azione e reazione.

\[
\left\{
\begin{array}{l}
\psi-kx_G+\xi_x=0,
\\[1mm]
-mg-Mg+\xi_y=0,
\\[1mm]
-R\xi_x=0,
\\[1mm]
mg\dfrac{l}{2}\sin\theta-l\psi\cos\theta=0.
\end{array}
\right.
\]

Quindi
\[
\left\{
\begin{array}{l}
\psi=kl\sin\theta,
\\[1mm]
\xi_y=mg+Mg,
\\[1mm]
mg\dfrac{l}{2}\sin\theta
-
l\cos\theta\,kl\sin\theta
=0.
\end{array}
\right.
\]

Da cui
\[
mg\sin\theta-2kl\sin\theta\cos\theta=0.
\]

Se
\[
\sin\theta=0,
\]
allora
\[
\theta_1=0,
\qquad
\theta_2=\pi.
\]

Se invece
\[
\sin\theta\neq 0,
\]
allora
\[
mg-2kl\cos\theta=0
\]
e quindi
\[
\cos\theta_{3,4}
=
\frac{mg}{2kl},
\qquad
\frac{mg}{2kl}\leq 1.
\]

Dal momento che
\[
\cos\theta_{3,4}>0,
\]
le configurazioni di equilibrio sono le seguenti.
\begin{center}
  \begin{tikzpicture}[scale=1, line cap=round, line join=round, >=Latex]

%======================
% Disegno 1
%======================
\begin{scope}[shift={(0,0)}]
    \node at (0,3.35) {$\theta_1=0$};

    % piano
    \draw[line width=0.9pt] (-1.7,0) -- (1.7,0);

    % disco
    \draw[line width=1pt] (0,1) circle (1);

    % asta verticale
    \draw[line width=1pt] (0,1) -- (0,2.85);

    % punto A e guida
    \fill (0,2.15) circle (1.2pt);
    \node[right] at (0,2.15) {$A$};
    \draw[line width=0.8pt] (-0.13,1.90) -- (-0.13,2.40);
    \draw[line width=0.8pt] ( 0.13,1.90) -- ( 0.13,2.40);
\end{scope}

%======================
% Disegno 2
%======================
\begin{scope}[shift={(5.6,0)}]
    \node at (0,3.35) {$\theta_2=\pi$};

    % piano
    \draw[line width=0.9pt] (-1.7,0) -- (1.7,0);

    % disco
    \draw[line width=1pt] (0,1) circle (1);

    % asta verticale verso il basso
    \draw[line width=1pt] (0,1) -- (0,-1.05);

    % punto A e guida
    \fill (0,-0.38) circle (1.2pt);
    \node[right] at (0,-0.38) {$A$};
    \draw[line width=0.8pt] (-0.13,-0.63) -- (-0.13,-0.13);
    \draw[line width=0.8pt] ( 0.13,-0.63) -- ( 0.13,-0.13);
\end{scope}

%======================
% Disegno 3
%======================
\begin{scope}[shift={(12.4,0)}]
    \node at (0,3.35) {$\theta_3,\theta_4$};

    % piano con freccia
    \draw[line width=0.9pt, ->] (-3.3,0) -- (4.2,0);

    % asse verticale centrale
    \draw[line width=1pt] (0,-1.4) -- (0,2.95);

    % dischi
    \draw[line width=1pt] (-1.7,1) circle (1);
    \draw[line width=1pt] ( 1.7,1) circle (1);

    % centri dei dischi
    \coordinate (GL) at (-1.7,1);
    \coordinate (GR) at ( 1.7,1);
    \coordinate (A)  at (0,2.18);

    % punto A e guida
    \fill (A) circle (1.2pt);
    \node[right] at (A) {$A$};
    \draw[line width=0.8pt] (-0.13,1.92) -- (-0.13,2.42);
    \draw[line width=0.8pt] ( 0.13,1.92) -- ( 0.13,2.42);

    % aste inclinate che partono dai centri delle circonferenze
    \draw[line width=1pt] (GL) -- (A);
    \draw[line width=1pt] (GR) -- (A);

    % molla orizzontale tra i due centri
    \draw[line width=1pt] (-1.7,1) -- (-1.25,1);

    \draw[line width=1pt, smooth, samples=100, domain=-1.25:-0.15]
        plot (\x,{1 + 0.07*sin((\x+1.25)/1.10*1080)});

    \fill (0,1) circle (1.2pt);

    \draw[line width=1pt, smooth, samples=100, domain=0.15:1.25]
        plot (\x,{1 + 0.07*sin((\x-0.15)/1.10*1080)});

    \draw[line width=1pt] (1.25,1) -- (1.7,1);
\end{scope}

\end{tikzpicture}
\end{center}

Troviamo le stesse posizioni con il principio dei lavori virtuali:
\[
-mg\hat{\jmath}\cdot \delta \underline{x}_B
-Mg\hat{\jmath}\cdot \delta \underline{x}_G
-kl\sin\theta\,\hat{\imath}\cdot \delta \underline{x}_G
=0
\qquad
\forall\,(\delta \underline{x}_B,\delta \underline{x}_G)
\text{ ammessi}.
\]

Passo alla coordinata libera:
\[
\underline{x}_B
=
\frac{l}{2}\sin\theta\,\hat{\imath}
+
\left(
R+\frac{l}{2}\cos\theta
\right)\hat{\jmath},
\qquad
\underline{x}_G
=
l\sin\theta\,\hat{\imath}
+
R\hat{\jmath}.
\]

Quindi
\[
\delta \underline{x}_B
=
\frac{l}{2}\cos\theta\,\delta\theta\,\hat{\imath}
-
\frac{l}{2}\sin\theta\,\delta\theta\,\hat{\jmath},
\qquad
\delta \underline{x}_G
=
l\cos\theta\,\delta\theta\,\hat{\imath}.
\]

Allora
\[
-mg\hat{\jmath}\cdot
\left(
-\frac{l}{2}\sin\theta\,\delta\theta\,\hat{\jmath}
\right)
-
Mg\hat{\jmath}\cdot
\left(
l\cos\theta\,\delta\theta\,\hat{\imath}
\right)
-
kl\sin\theta\,\hat{\imath}\cdot
\left(
l\cos\theta\,\delta\theta\,\hat{\imath}
\right)
=0.
\]

Da cui
\[
mg\frac{l}{2}\sin\theta\,\delta\theta
-
kl\sin\theta\,l\cos\theta\,\delta\theta
=0
\qquad
\forall\,\delta\theta.
\]

Pertanto
\[
mg\frac{l}{2}\sin\theta
-
kl\sin\theta\,l\cos\theta
=0.
\]

Equivalentemente,
\[
mg\sin\theta
-
2kl\sin\theta\cos\theta
=0.
\]

Quindi
\[
\theta_1=0,
\qquad
\theta_2=\pi,
\qquad
\cos\theta_{3,4}
=
\frac{mg}{2kl},
\qquad
\frac{mg}{2kl}\leq 1.
\]
\end{example}

\vspace{60pt}
\begin{example}{}
  Studiamo la statica del seguente sistema.
  \begin{center}
    \begin{tikzpicture}[
    x=1cm,
    y=1cm,
    >=Latex,
    line cap=round,
    line join=round,
    every node/.style={font=\small}
]

% Coordinate principali
\coordinate (O) at (0,0);
\coordinate (A) at (2.70,0);
\coordinate (C) at (6.10,0);
\coordinate (B) at (6.10,-2.35);
\coordinate (G) at ($(A)!0.47!(B)$);

% Riferimento a sinistra
\draw[blue!70!black, line width=1.4pt, ->] (O) -- ++(0,1.10);
\draw[blue!70!black, line width=1.4pt, ->] (O) -- ++(0.85,0);
\node[blue!70!black, left]  at ($(O)+(0,0.60)$) {$\underline{\jmath}$};
\node[blue!70!black, below] at ($(O)+(0.42,0)$) {$\underline{\imath}$};

\draw[blue!70!black, line width=1.2pt] ($(O)+(-0.55,-0.45)$) circle (0.20);
\fill[blue!70!black] ($(O)+(-0.55,-0.45)$) circle (0.035);
\node[blue!70!black, left]  at ($(O)$) {$O$};
\node[blue!70!black, below] at ($(O)+(-0.55,-0.72)$) {$\underline{k}$};

% Molla sinistra attaccata al riferimento
\draw[yellow!80!orange!90!black, line width=1.4pt] (O) -- (0.22,0);
\draw[yellow!80!orange!90!black, line width=1.4pt, smooth, samples=140, domain=0.22:2.48]
plot (\x,{0.16*sin(deg((\x-0.22)*15.5))});
\draw[yellow!80!orange!90!black, line width=1.4pt] (2.48,0) -- (A);
\node at (1.00,0.55) {$K$};

% Tratto orizzontale superiore
\draw[black, line width=1.2pt] (0,0) -- (7.05,0);

% Asta inclinata
\draw[black, line width=1.2pt] (A) -- (B);

% Molla destra verticale
\draw[yellow!80!orange!90!black, line width=1.4pt] (C) -- ++(0,-0.22);
\draw[yellow!80!orange!90!black, line width=1.4pt, smooth, samples=140, domain=-0.22:-2.13]
plot ({6.10 + 0.10*sin(deg((\x+0.22)*16.5))}, \x);
\draw[yellow!80!orange!90!black, line width=1.4pt] (6.10,-2.13) -- (B);
\node[right] at ($(C)+(0.45,-1.15)$) {$K$};

% Punti e nomi
\node[below left]  at (A) {$A$};
\node[above]       at (C) {$C$};
\node[below right] at (B) {$B$};
\node[below]       at ($(B)+(-0.45,-0.10)$) {$m,l$};

\fill[black] (G) circle (1.6pt);
\node[right] at ($(G)+(0.05,0.05)$) {$G$};

% Angolo theta
\draw[black, line width=1.0pt]
($(A)+(0.65,0)$) arc[start angle=0,end angle=-38,radius=0.65];
\node at ($(A)+(0.82,-0.40)$) {$\theta$};

% Freccia rossa F
\draw[red!80!black, line width=1.5pt, ->] ($(A)+(0.05,0)$) -- ++(1.05,0);
\node[red!80!black, above] at ($(A)+(0.55,0.38)$) {$\underline{F}$};

% Freccia verde psi in A
\draw[green!55!black, line width=1.4pt, ->] ($(A)+(0,-0.05)$) -- ++(0,-0.60);
\node[green!55!black, right] at ($(A)+(0.10,-0.47)$) {$\underline{\psi}$};

% Forza elastica sinistra
\draw[red!80!black, line width=1.5pt, ->] ($(A)+(-0.25,-0.55)$) -- ++(-1.10,0);
\node[red!80!black, below] at ($(A)+(-0.80,-0.55)$) {$-k\,\underline{x}_A$};

% Peso in G
\draw[red!80!black, line width=1.5pt, ->] (G) -- ++(0,-1.00);
\node[red!80!black, left] at ($(G)+(-0.10,-0.68)$) {$m\underline{g}$};

% Forza elastica destra
\draw[red!80!black, line width=1.5pt, ->] ($(C)+(1.20,-2.00)$) -- ++(0,1.10);
\node[red!80!black, right] at ($(C)+(1.40,-1.45)$)
{$K\big(\underline{x}_C-\underline{x}_B\big)$};

\end{tikzpicture}
  \end{center}
  Imposto un riferimento fisso
\[
\{O,\hat{\imath},\hat{\jmath},\hat{k}\}.
\]
L'asta di lunghezza \(l\) e massa \(m\) è sottoposta al vincolo
\[
y_A^{\mathrm{asta}}=0
\]
ed ha quindi \(2\) gradi di libertà. Scelgo come coordinate libere
\[
x_A,\theta.
\]
Studiamo le sue equazioni della statica:
\[
\left\{
\begin{array}{l}
\underline{F}
-k\underline{x}_A
+\underline{\psi}
+m\underline{g}
+k\left(\underline{x}_C-\underline{x}_B\right)
=
\underline{0},
\\[2mm]
\left(\underline{x}_G-\underline{x}_A\right)\wedge m\underline{g}
+
\left(\underline{x}_B-\underline{x}_A\right)\wedge
k\left(\underline{x}_C-\underline{x}_B\right)
=
\underline{0}
\qquad
\text{polo } A.
\end{array}
\right.
\]

Quindi
\[
\left\{
\begin{array}{l}
F\hat{\imath}
-kx_A\hat{\imath}
-\psi\hat{\jmath}
-mg\hat{\jmath}
+kl\sin\theta\,\hat{\jmath}
=
\underline{0},
\\[2mm]
\dfrac{l}{2}
\left(
\cos\theta\,\hat{\imath}
-
\sin\theta\,\hat{\jmath}
\right)
\wedge
\left(
-mg\hat{\jmath}
\right)
+
l
\left(
\cos\theta\,\hat{\imath}
-
\sin\theta\,\hat{\jmath}
\right)
\wedge
kl\sin\theta\,\hat{\jmath}
=
\underline{0}.
\end{array}
\right.
\]

Passando alle componenti:
\[
\left\{
\begin{array}{l}
F-kx_A=0,
\\[1mm]
-\psi-mg+kl\sin\theta=0,
\\[1mm]
-\dfrac{l}{2}\cos\theta\,mg
+
l\cos\theta\,kl\sin\theta
=0.
\end{array}
\right.
\]

Dalla prima equazione:
\[
x_A=\frac{F}{k}.
\]

Dalla terza equazione:
\[
-\frac{l}{2}\cos\theta\,mg
+
kl^2\sin\theta\cos\theta
=0,
\]
cioè
\[
2kl\sin\theta\cos\theta
-
mg\cos\theta
=0.
\]

Se
\[
\cos\theta=0,
\]
allora
\[
\theta_1=\frac{\pi}{2},
\qquad
\theta_2=\frac{3}{2}\pi.
\]

Se invece
\[
\cos\theta\neq 0,
\]
allora
\[
2kl\sin\theta=mg,
\]
da cui
\[
\sin\theta=\frac{mg}{2kl}.
\]

Quindi
\[
\sin\theta_{3,4}
=
\frac{mg}{2kl},
\qquad
\frac{mg}{2kl}\leq 1.
\]

Facciamolo con il principio dei lavori virtuali:
\[
-k\underline{x}_A\cdot \delta\underline{x}_A
+
\underline{F}\cdot \delta\underline{x}_A
+
m\underline{g}\cdot \delta\underline{x}_G
+
k\left(\underline{x}_C-\underline{x}_B\right)\cdot
\delta\underline{x}_B
=
0.
\]

Scrivo gli spostamenti virtuali in coordinate libere:
\[
\underline{x}_A
=
x_A\hat{\imath}
\qquad \Longrightarrow \qquad
\delta\underline{x}_A
=
\delta x_A\hat{\imath}.
\]

\[
\underline{x}_G
=
\underline{x}_G-\underline{x}_A+\underline{x}_A
=
\frac{l}{2}
\left(
\cos\theta\,\hat{\imath}
-
\sin\theta\,\hat{\jmath}
\right)
+
x_A\hat{\imath},
\]
quindi
\[
\delta\underline{x}_G
=
\delta x_A\hat{\imath}
-
\frac{l}{2}\sin\theta\,\delta\theta\,\hat{\imath}
-
\frac{l}{2}\cos\theta\,\delta\theta\,\hat{\jmath}.
\]

Analogamente,
\[
\underline{x}_B
=
\underline{x}_B-\underline{x}_A+\underline{x}_A
=
l
\left(
\cos\theta\,\hat{\imath}
-
\sin\theta\,\hat{\jmath}
\right)
+
x_A\hat{\imath},
\]
quindi
\[
\delta\underline{x}_B
=
\delta x_A\hat{\imath}
-
l\sin\theta\,\delta\theta\,\hat{\imath}
-
l\cos\theta\,\delta\theta\,\hat{\jmath}.
\]

Sostituendo nel principio dei lavori virtuali:
\[
-Kx_A\hat{\imath}\cdot \delta x_A\hat{\imath}
+
F\hat{\imath}\cdot \delta x_A\hat{\imath}
-
mg\hat{\jmath}\cdot
\left(
\delta x_A\hat{\imath}
-
\frac{l}{2}\sin\theta\,\delta\theta\,\hat{\imath}
-
\frac{l}{2}\cos\theta\,\delta\theta\,\hat{\jmath}
\right)
\]
\[
+
Kl\sin\theta\,\hat{\jmath}\cdot
\left(
\delta x_A\hat{\imath}
-
l\sin\theta\,\delta\theta\,\hat{\imath}
-
l\cos\theta\,\delta\theta\,\hat{\jmath}
\right)
=
0.
\]

Da cui
\[
\left\{
\begin{array}{l}
-Kx_A\,\delta x_A
+
F\,\delta x_A
=
0,
\\[2mm]
mg\dfrac{l}{2}\cos\theta\,\delta\theta
-
Kl\sin\theta\,l\cos\theta\,\delta\theta
=
0,
\end{array}
\right.
\qquad
\forall\,(\delta x_A,\delta\theta).
\]

Quindi
\[
\left\{
\begin{array}{l}
F=Kx_A,
\\[2mm]
mg\dfrac{l}{2}\cos\theta
-
Kl\sin\theta\,l\cos\theta
=
0.
\end{array}
\right.
\]

Equivalentemente,
\[
\left\{
\begin{array}{l}
x_A=\dfrac{F}{K},
\\[3mm]
mg\cos\theta
-
2Kl\sin\theta\cos\theta
=
0.
\end{array}
\right.
\]

E arriviamo alla stessa conclusione di prima.
\end{example}
